\documentclass[10pt]{article}
\usepackage[utf8]{inputenc}
\usepackage{url}
\usepackage{hyperref}
\usepackage{amsmath}
\usepackage{amsfonts}
\usepackage{amssymb}
\usepackage{graphicx}
\usepackage{float}
\usepackage{lipsum}
\usepackage{multicol}
\usepackage{xcolor}
\usepackage{tabularx}
\usepackage{natbib}
\usepackage[font=small]{caption}
\addtolength{\abovecaptionskip}{-3mm}
\addtolength{\textfloatsep}{-5mm}
\setlength\columnsep{20pt}

\bibliographystyle{plainnat}
\setcitestyle{numbers,square}

\usepackage[a4paper,left=1.50cm, right=1.50cm, top=1.50cm, bottom=1.50cm]{geometry}


\author{}

\title{Project Proposal -- CIS*4900}

\begin{document}
	
	\begin{center}
		{\Large \textbf{Project Proposal: Bayesian Estimation}}\\
		\vspace{1em}
		\textit{\text{Stella Rovazzi} $\cdot$ \text{Ramez Elprince} $\cdot$ \text{Brenden Doogan}}\\
		\vspace{1em}
		\textit{University of Guelph}
	\end{center}
	

	\begin{center}
		\rule{150mm}{0.2mm}
	\end{center}		

	
% \begin{multicols*}{2}
 
\section*{Project Overview} 

\subsection*{Objectives}

\begin{itemize}
\item bleh \cite{ripley_mass_2009} 
\end{itemize}

\subsection{Background and Literature}

\section{Proposed Computational Strategy}
\subsection*{Simple Example}
\subsubsection*{Goal}
We will start with a simple example to ensure we can correctly converge to some posterior.
Alongside this, we will ensure we can calculate (and infer) posterior summary statistics,
suitability of the model, monte-carlo standard error, and rate of convergence.

\subsubsection*{Methodology}
This follows the example C.3 regarding the normals in Bayesain Data Analysis (BDA) \cite{gelman_bayesian_2014}.

\begin{enumerate}
	\item Generate bivariate normal data
	\item Derive and make functions for the marginal distributions (assuming normal prior and conditional)
	\item Run the posterior simulations with Gibbs Sampling
	\item Find the parameters through the samples
	\item Repeat for various realistic n
\end{enumerate}

\subsection*{Harder Example}
\subsubsection*{Goal}
Apply what was learned in the simple example to real data. We will be able to try and compare the our methodology. 
\subsubsection*{Methodology}
This follows the 1.6 example in BDA \cite{gelman_bayesian_2014} but we turn it computational. 

\begin{itemize}
	\item Derive the full conditional, then the marginal conditionals using the distributions provided in text
	\item Gibbs sample for the posterior
	\item Explore the results
	\item Repeat for various realistic n
\end{itemize}

\subsection*{Extend to Hierarchical Models}


\subsection*{Quantities of Interest}

\begin{itemize}
	\item Estimated mean, variance/standard deviation, credible intervals (informs us of the posterior model)
	\item Gelman-Rubin Statistic (Convergence Assessment) 
	\item Rate of Convergence (computational efficiency)
\end{itemize}

\subsection*{Required Libraries}
\begin{itemize}
	\item MASS  for the bivariate normal generation (or basic R) \cite{ripley_mass_2009}
	\item STAN for the actual MCMC/Gibs sampling – we may some examples by hand \cite{guo_rstan_2015}
\end{itemize}

\section*{Division of Responsibility}


\section*{Timeline}

**We can chat about GANTT charts on tuesday or smth, just so we can layout realistic dates with buffer room**.
\section*{Expected Outcomes}

bleh

% \end{multicols*}{2}

\bibliography{bb-ds-fellowship}
	
\end{document}